\section{Introduction}
\label{sec:introduction}

Repeated computation can concentrate representations into a small subspace. This observation appears in graph oversmoothing \citep{oono2020exponential,roth2024rank}, attention rank loss \citep{pure_attention_rank}, and dimensional collapse in representation learning \citep{jing2022dimensional}. Its significance depends on what the representation must preserve. A compact state may encode every required variable; a high-dimensional state may consist largely of irrelevant fluctuations. In recurrent models, a further ambiguity arises: local contraction can coexist with useful memory when different cues select different stable responses \citep{pals2024-phase-locked-memory}. A geometric observation alone cannot distinguish efficient compression from destructive forgetting.

We consider a concrete question: after a recurrent model has acquired task-required information, does a periodic distractor body cause more loss than a matched non-periodic body of the same length? The network's weights, written values, query, and number of updates remain fixed across the comparison. The intervention changes motif ordering while matching symbol counts, bigram counts, run lengths, endpoints, and the final motif. This isolates a specified ordering contrast, rather than a universal effect of ``periodicity'': higher-order input statistics necessarily remain different.

Answering the question requires distinguishing three objects. A covariance spectrum describes the geometry of a selected ensemble. A tangent product describes propagation of infinitesimal state perturbations along a particular trajectory. Native behavior and controlled decoding describe different forms of access to task information. Their relationships require assumptions. Isotropic attenuation preserves normalized covariance dimension while reducing signal relative to fixed noise. Nuisance compression lowers dimension while preserving the target. Separate stable basins retain categorical labels despite contracting local spectra. An arm-dependent readout offset changes error even when the cue-dependent state component is identical. \Cref{fig:controls} makes these distinctions concrete.

We contribute a compact set of analytic controls, a matched-body measurement design, and an empirical audit that demonstrates why verified native competence must precede a harmful-retention claim. The controls assemble established dynamical and information-theoretic ideas into explicit falsifiers: fixed affine and commuting switched-affine systems cannot support particular ordering explanations; local stability does not determine global cue retention; and normalized dimension does not determine noise-resolved information. We use \emph{diagnostic non-identifiability} in this limited sense: a diagnostic value need not determine whether the required information survives. The individual ingredients are not claimed as new dynamical theorems.

The audit covers four trained model/task combinations: GRU and LSTM, each with categorical and continuous targets. Every selected checkpoint fails short-horizon competence. Consequently, their later errors cannot establish erasure of successfully acquired information. This negative outcome is useful as a methodological case study, but it supplies neither a periodicity effect nor evidence of general recurrent collapse. The paper therefore centers the inference problem and its controls, with the remaining empirical question stated explicitly.
